\documentclass[12pt]{article}
\usepackage[utf8]{inputenc}
\usepackage[T1]{fontenc}
\usepackage{hyperref}
\title{Wawa Language Resources for K-12 Curriculum Development}
\author{Compiled by Claude Sonnet 5.5 (claude-sonnet-5-5) for Andr\'as Kornai}
\date{}
\begin{document}
\maketitle

\section{Introduction}
Wawa (ISO 639-3 www; Glottocode wawa1246) is a threatened Mambiloid language of the Niger-Congo family. It is spoken in the Bankim area of the Mayo-Banyo Division, Adamawa Region, Cameroon, and extends slightly into Nigeria. Wikipedia gives about 3,000 speakers in three main dialects, while Joshua Project gives a higher figure. Griffiths and Robson (2010), as summarised in search results, report that children still acquire the language and use it daily. They also report that it is unwritten and that children are taught to read the Koran. They flag Fulfulde as a potential threat.

\section{Primary Educational Resources}
No primers, textbooks or school programmes for Wawa appeared in the searches. The language has no established orthography in the sources seen.

\section{Linguistic and Reference Materials}
\begin{itemize}
\item Marieke Martin (2012), \emph{A grammar of Wawa: an endangered language of Cameroon}. It is cited as the main reference. I saw no direct URL for it, so it is not listed in the JSON.
\item Wikipedia article on Wawa. \url{https://en.wikipedia.org/wiki/Wawa_language}
\item Endangered Languages Project entry for Wawa. \url{https://endangeredlanguages.com/lang/6105?hl=en}
\item Numerals page for Wawa from the MPI-EVA Chan numerals database. \url{https://lingweb.eva.mpg.de/channumerals/Wawa.htm}
\end{itemize}

\section{Religious and Community Resources}
Joshua Project has a people-group profile for the Wawa in Cameroon. \url{https://www.joshuaproject.net/index.php/people_groups/15895} No scripture translation or community organisation was found in the searches.

\section{Digital and Audio Resources}
The Endangered Languages Project page has a language sample page. \url{https://www.endangeredlanguages.com/lang/6105/samples/10989?hl=fr} I did not verify its content. No other audio, video or software was found.

\section{Curriculum Development Context}
Wawa is unwritten and, as far as the searches show, not taught in schools. A curriculum would first need a community-agreed orthography, and it would need to build on Martin's grammar. The community's own priorities should lead. Teaching materials are the main gap.
\end{document}